% Adaptado por Andre
% ============================================================
% PLANTILLA JAES -- ESTILO DE CITACIÓN: IEEE NUMÉRICO
% ============================================================
% Citas como [1], [2], [3] en orden de aparición en el texto.
% Bibliografía ordenada por orden de cita (NO alfabético).
% Usa: jaes.bst (el estilo "nativo" de esta plantilla, NO es de
%      TeX Live/MiKTeX -- viene incluido en esta misma carpeta).
%
% Úsala si: el profesor/revista no exige un estilo en particular,
% o pide explícitamente formato IEEE / numérico.
%
% ---- Entorno de referencia de esta plantilla ----
%   SO:           Windows 11
%   Editor:       VS Code + extensión LaTeX Workshop
%   Motor:        latexmk (vía Strawberry Perl)
%   Distribución: MiKTeX, con auto-instalación de paquetes activa
%
% Archivos que deben estar en esta misma carpeta:
%   - plantilla_ieee.tex  (este archivo)
%   - jaes.cls            (clase del documento, define el formato JAES)
%   - jaes.bst            (estilo bibliográfico numérico)
%   - jaes.bib            (tu base de datos de referencias)
%
% Compilación con LaTeX Workshop + latexmk:
%   No necesitas correr pdflatex/bibtex a mano. Con latexmk configurado
%   (perfil por defecto de LaTeX Workshop ya lo trae), basta con guardar
%   el archivo o pulsar el botón "Build LaTeX project" (ícono ▶ verde).
%   latexmk detecta solo que hay un \bibliography{} + \bibliographystyle{}
%   y corre bibtex en el orden correcto las veces que haga falta.
%   La primera vez que MiKTeX vea un paquete que falte (poco probable
%   en esta plantilla, ya que todo es estándar), te preguntará si
%   quieres instalarlo: acepta, y MiKTeX lo descarga e instala solo.
% ============================================================

\documentclass[fleqn, final]{jaes}

% ---- Metadatos del pie de página ----
\jyear{2026}                          % Año que aparece en el pie de página
\jmonth{Junio}                        % Mes que aparece en el pie de página
\jname{Nombre del Curso -- Código}     % Texto del pie (por defecto: "Proyecto de Investigación")

% ---- Paquetes adicionales ----
% Todos estos son paquetes ESTÁNDAR (MiKTeX los autoinstala la primera
% vez que se usan, sin que tengas que hacer nada manual).
\usepackage[utf8]{inputenc}
\usepackage{amsmath}\setlength{\mathindent}{10pt}
\usepackage{bm}

% cleveref: permite \cref{} / \Cref{} para referencias automáticas
% ("Sec. 3", "Tabla 2", etc.) sin escribir el número a mano.
\usepackage[capitalise]{cleveref}
  \crefformat{section}{\textsc{Sec.~#2#1#3}}
  \Crefformat{section}{\textsc{Sec.~#2#1#3}}
  \crefformat{figure}{{Fig.~#2#1#3}}
  \Crefformat{figure}{{Fig.~#2#1#3}}
  \crefformat{table}{{Tabla~#2#1#3}}
  \Crefformat{table}{{Tabla~#2#1#3}}
  \crefname{section}{Sec.}{Secs.}
  \Crefname{section}{Sec.}{Secs.}
  \crefname{table}{Tabla}{Tablas}
  \Crefname{table}{Tabla}{Tablas}
  \crefname{figure}{Fig.}{Figs.}
  \Crefname{figure}{Fig.}{Figs.}
% Fix: cleveref dice "and" en vez de "y" al combinar 2+ referencias
% (\cref{sec:a,sec:b}) si el documento no usa el paquete babel.
% Sin este bloque, saldría "(Sec. 2 and 5)" en inglés.
\providecommand{\creflastconjunction}{}
\renewcommand{\creflastconjunction}{ y }
\providecommand{\crefpairconjunction}{}
\renewcommand{\crefpairconjunction}{ y }
\providecommand{\crefmiddleconjunction}{}
\renewcommand{\crefmiddleconjunction}{, }

% Paquetes para tablas anchas / con texto que necesita ajustarse de línea
% (muy recomendable si vas a tener tablas con párrafos largos en celdas).
\usepackage{array}
\usepackage{tabularx}
\usepackage{longtable}
\usepackage{multirow}
\usepackage{ragged2e}
\newcolumntype{Y}{>{\RaggedRight\arraybackslash}X}  % columna de ancho fijo con ajuste de línea

% Localización al español (jaes.cls viene en inglés por defecto):
\renewcommand{\tablename}{Tabla}      % por defecto decía "Table"
% \bibname (NO \refname) es el comando que jaes.cls usa internamente
% para titular la sección de bibliografía (\section{\bibname} dentro
% del entorno thebibliography). Por defecto vale "REFERENCES".
\renewcommand{\bibname}{Referencias}

% Anula la biografía automática de autores (\biography), que tiene un
% bug conocido en esta clase: no imprime el texto biográfico, solo
% fotos y nombres. Ver la sección "Sobre el Equipo de Trabajo" más abajo
% para la alternativa manual que SÍ funciona.
\def\printbiography{}

% ---- Si tu entrega NO debe decir "BORRADOR" en cada página ----
% jaes.cls imprime la palabra BORRADOR en el encabezado de todas las
% páginas de forma fija (no es una opción togglable desde el .tex).
% Para quitarlo hay que editar tu copia de jaes.cls, buscando estas
% dos líneas y dejando vacío el contenido de \rhfont:
%
%   \def\@evenhead{\hbox to \textwidth{\hfill\rhfont BORRADOR}}%
%   \def\@oddhead{\hbox to \textwidth{\rhfont BORRADOR\hfill}}%
%
% (Si tu archivo aún es un borrador en revisión, puedes dejarlo así).

\begin{document}

% Encabezado de cada página: {Lista de autores}{Título corto del trabajo}
\markboth{APELLIDO1, APELLIDO2, APELLIDO3}{Título Corto del Trabajo}

% ---------------- Título ----------------
\title{Título Completo del Trabajo}

% ---------------- Autores ----------------
% Un \author{...} por persona, separados por comas, todos dentro de
% un único \authorgroup{...}. El \textsuperscript{1} enlaza cada autor
% con su afiliación correspondiente en \affil{}.
% IMPORTANTE: \email{} y \affil{} deben ir DENTRO del bloque
% \authorgroup{...} (antes de la llave de cierre), no después. Si se
% colocan fuera, su formato (cursiva, tamaño de fuente) se filtra al
% resto del documento porque \authorgroup solo GUARDA su contenido
% para imprimirlo más tarde con \maketitle; lo que esté fuera de su
% llave de cierre se ejecuta de inmediato, en un punto del documento
% donde el diseño de página todavía no está listo para recibirlo.
\authorgroup{
\author{NOMBRE APELLIDO\textsuperscript{1}\orcidlink{0000-0000-0000-0000}},
\author{NOMBRE APELLIDO 2\textsuperscript{1}},
\author{NOMBRE APELLIDO 3\textsuperscript{1}}
% \email{...} es opcional: úsalo para correos, códigos de alumno, etc.
\email{(código1) \quad (código2) \quad (código3)}
\affil{\textsuperscript{1}Universidad, Facultad, Carrera, Ciudad, País}
}

% ---------------- Resumen ----------------
% El abstract va en un solo párrafo, sin saltos de línea en blanco.
% Longitud orientativa: 150-250 palabras. Debe poder leerse de forma
% independiente del resto del documento (qué se hizo, cómo, y qué se
% encontró), sin citar tablas/figuras por número.
\abstract{%
Texto del resumen aquí. Describe el problema u objetivo, la metodología
utilizada, los principales hallazgos y la conclusión o recomendación
central del trabajo, en un solo párrafo corrido.}
\maketitle

% ---------------- Palabras clave ----------------
\begin{keywords}
palabra clave 1, palabra clave 2, palabra clave 3, palabra clave 4
\end{keywords}

% ============================================================
% CUERPO DEL DOCUMENTO
% Reemplaza estos títulos de sección por los que pida tu rúbrica/curso.
% \label{sec:...} permite referenciar la sección con \cref{sec:...}
% desde cualquier otra parte del documento.
% ============================================================

\section{Introducción / Diagnóstico}\label{sec:introduccion}

Texto de la sección. Usa \verb|\subsection{...}| para subdividir si
la sección es larga.

\subsection{Subsección de ejemplo}

Para citar una fuente del archivo .bib, usa \verb|\cite{clave}|, que en
este estilo (IEEE) se muestra como un número entre corchetes, ej.
\cite{ejemplo_clave}.

% ---- Ejemplo de tabla simple (una columna) ----
\begin{table}[tb]
\tabcolsep8.1pt
\centering
\caption{Título descriptivo de la tabla.}
\label{table:ejemplo}
{%
\begin{tabular}{@{}lcc@{}}\toprule
Columna 1 & Columna 2 & Columna 3\\\colrule
Fila 1 & dato & dato\\
Fila 2 & dato & dato\\\botrule
\end{tabular}}
\end{table}

Se referencia así: ver \cref{table:ejemplo}.

% ---- Ejemplo de tabla ancha (a dos columnas, usa tabularx) ----
% Usar table* (no table) para que ocupe el ancho completo de la página.
% La columna tipo Y permite que el texto largo haga salto de línea
% automático en vez de desbordarse fuera de la página.
\begin{table*}[tb]
\centering
\caption{Título descriptivo de la tabla ancha.}
\label{table:ejemplo_ancho}
{\small
\begin{tabularx}{\textwidth}{@{}lYc@{}}\toprule
Columna corta & Columna de texto largo (ajusta línea) & Columna corta\\\colrule
Dato & Texto largo que se ajustará automáticamente dentro del ancho de columna asignado, sin desbordarse. & Dato\\\botrule
\end{tabularx}}
\end{table*}

\section{Metodología}\label{sec:metodologia}

Texto de la sección.

\section{Resultados / Análisis}\label{sec:resultados}

Texto de la sección.

\section{Conclusiones}\label{sec:conclusiones}

Texto de la sección.

% ============================================================
% Bibliografía -- ESTILO IEEE NUMÉRICO
% ============================================================
\bibliography{jaes}            % nombre del .bib SIN extensión
\bibliographystyle{jaes}       % estilo numérico nativo de esta plantilla

% ============================================================
% Equipo de trabajo (alternativa manual a \biography, que no
% imprime el texto biográfico correctamente en esta clase)
% ============================================================
\clearpage
\markboth{Equipo de Trabajo}{Equipo de Trabajo}
\section*{Sobre el Equipo de Trabajo}

\begin{description}
\item[Nombre Apellido (código)] Rol o aporte específico al trabajo.
\item[Nombre Apellido 2 (código)] Rol o aporte específico al trabajo.
\item[Nombre Apellido 3 (código)] Rol o aporte específico al trabajo.
\end{description}

\noindent Texto de cierre opcional: institución, curso, docente, etc.

\end{document}
